\chapter{Background and Related Work}\label{chap:background}


% =====================================================================
\section{Autoregressive Decoding and the KV-Cache Bottleneck}\label{sec:kv-bottleneck}
% =====================================================================

Decoder transformers~\cite{vaswani2017attention} generate text
autoregressively: at decode step $t$, the model conditions on every
previously emitted token and produces a single new token.  The
attention mechanism at each layer computes, for the current
query~$q_t \in \mathbb{R}^{d}$ and the cumulative keys
$K_{\le t} \in \mathbb{R}^{t \times d}$ and values
$V_{\le t} \in \mathbb{R}^{t \times d}$,
\begin{equation}\label{eq:attn}
  \mathrm{Attn}(q_t, K_{\le t}, V_{\le t})
  \;=\;
  \mathrm{softmax}\!\bigl(q_t K_{\le t}^{\top} / \sqrt{d}\bigr)\,V_{\le t}.
\end{equation}
To avoid recomputing the past keys and values at every step, modern
inference runtimes maintain a \emph{KV cache} that grows by one row
per decode step.  Each step then performs an inner product between
$q_t$ and the entire cached $K_{\le t}$, followed by a softmax-weighted
read of $V_{\le t}$.

The cost of Eq.~\eqref{eq:attn} per decode step is dominated by the
memory traffic of reading the cache.  At context length $L$, the per-step
attention reads $O(L \cdot d)$ entries of $K$ and $V$ from GPU global
memory, and the arithmetic intensity of the operation (a single small
matrix--vector product per head) is far below the compute-to-memory
ratio of modern accelerators.  As a result, decode latency scales
approximately linearly with $L$ and is bottlenecked by HBM bandwidth
rather than by FLOPs.  At $L = 128$K with an 8B-parameter model in
\texttt{bfloat16}, the KV cache alone can occupy more than half of the
device memory budget, and the per-step attention can take longer than
all other layer operations combined.

Two structural observations sit behind every sparse-attention
proposal of the past two years.  First, even at long context the
dense softmax in Eq.~\eqref{eq:attn} concentrates most of its mass
on a small subset of past tokens; the rest contributes negligible
energy and could in principle be skipped without changing the output.
Second, prefill --- the one-pass attention over the prompt that
populates the cache --- is compute-bound rather than memory-bound.
Concretely, prefill on an $L$-token prompt computes a $Q[L \times d]
\cdot K^{\top}[d \times L]$ matrix--matrix product (arithmetic
intensity $O(L)$ FLOP per byte loaded), whereas decode computes a
$q[1 \times d] \cdot K^{\top}[d \times L]$ matrix--vector product
(arithmetic intensity $O(1)$, independent of $L$).  For any $L$ on
the order of a thousand or more, prefill is therefore in the
compute-bound regime and modifying it offers little speedup, while
decode is squarely memory-bound.  Sparse-attention
systems --- DCT-Page included --- therefore typically leave prefill
unchanged and intervene only at decode time, replacing the full read
of $K_{\le t}, V_{\le t}$ with some compressed or filtered surrogate.


% =====================================================================
\section{Sparse and Compressed KV Attention}\label{sec:related-sparse}
% =====================================================================

The methods DCT-Page is compared against differ along two structural
axes: (i) how they \emph{group} the KV cache into selectable units
--- fixed-size pages or learned clusters --- and (ii) how they
\emph{score} each unit against the current query.  The first three
methods below share the page-based paradigm with DCT-Page; Multipole
sits in the cluster-based paradigm and is discussed last.

\subsection{Quest}\label{sec:rel-quest}

Quest~\cite{tang2024quest} partitions the KV cache into fixed-size
pages and, for each page, stores a pair of per-channel summary
statistics: the elementwise minimum and maximum of the keys inside
the page.  At decode time the query is scored against every page by an
\emph{approximate-max inner product}.  For each channel
$c \in \{1, \dots, d\}$, the per-channel contribution is taken to be
the larger of $q[c] \cdot \min_p[c]$ and $q[c] \cdot \max_p[c]$ ---
since the true key value $k[c]$ for any token inside the page lies in
$[\min_p[c], \max_p[c]]$, this is a tight upper bound on
$q[c] \cdot k[c]$ --- achievable when some token in the page attains
the channel extreme --- and selecting between $\min_p[c]$ and
$\max_p[c]$ requires only the per-channel sign of $q[c]$ ($\max_p[c]$
if $q[c] \ge 0$, $\min_p[c]$ otherwise).
The page score is the sum of these per-channel upper bounds across
the $d$ channels, which is itself an upper bound on
$\max_{k \in \text{page}} q^{\top} k$.  This bound is cheap (two
scalars per channel) and position-agnostic.  Quest then attends to the
top-$k$ pages by this score and ignores the rest.

The per-page metadata is small --- two scalars per channel rather
than a learned proxy --- and the kernel cost of scoring is low, which
makes Quest fast at moderate context length.  However, the min/max
abstraction discards \emph{intra-page} structure entirely: it cannot
distinguish a page in which a few channels happen to span a wide
range from a page in which one specific token has a large inner
product with the query.  Section~\ref{sec:quest-collapse} returns to
this point empirically.

\subsection{SeerAttention-R}\label{sec:rel-seer}

SeerAttention-R~\cite{gao2025seerattention} adds a small learned
\emph{gate} to each decoder layer that, at decode time, predicts
per-block importance scores from the current query.  The gate is
trained on long-context data with a supervision signal derived from
the dense attention distribution, so that its block scores correlate
with the softmax mass that the dense model would have placed on each
block.  At inference, the gates are evaluated at every layer and the
top-$k$ blocks are kept under a fixed token budget.

The learned-gate approach can be very accurate at a given budget
because it is explicitly optimized for the selection task, but it
imposes two practical constraints.  First, a model-specific
checkpoint must be released for every base model; in our experiments
we use the Qwen3-family gates released by the authors, since no
equivalent Llama-3 gates are available, and SeerAttention-R is
therefore reported on Qwen3-8B only.  Second, the gate adds a small
but non-zero per-layer parameter count and forward-time cost, which
slightly reduces its ceiling on raw decode throughput.

\subsection{InfLLM}\label{sec:rel-infllm}

InfLLM~\cite{xiao2024infllm} approaches long-context decoding from a
retrieval perspective.  The past KV cache is broken into blocks; each
block is summarized by a small set of \emph{representative tokens}
chosen by an upstream-defined importance heuristic; and at decode
time the current query is matched against every block's
representatives to produce a relevance score.  InfLLM additionally
keeps an initial sink region (the first $n_\text{init}$ tokens) and a
local sliding window, attending in full to both regardless of the
retrieval result.  Only the top-ranked blocks outside the sink and
local zones are read at full precision.

The retrieval framing makes InfLLM agnostic to context length in
principle, but it pays for this flexibility in implementation
complexity: the block representatives must be selected and indexed
carefully, and the published code path currently targets the Llama-3
family only.  We accordingly report InfLLM on Llama-3.1-8B-Instruct
alone.

\subsection{Multipole Attention}\label{sec:rel-multipole}

Multipole Attention~\cite{hooper2025multipole} departs from the
fixed-page paradigm and instead clusters keys with hierarchical
$k$-means.  At decode time the query attends through cluster
centroids: each cluster contributes a single representative key/value
pair, and clusters whose centroids score highly are then expanded to
their constituent keys for the final attention.  This is the
attention analogue of the fast-multipole-method approximation used in
$n$-body simulation.

Cluster-based selection has two strengths and one structural
weakness.  Its strengths are that the cluster representatives carry
within-page structure (a centroid is a positionless but
distribution-aware summary) and that the hierarchical structure
allows variable-width effective receptive fields per layer.
Empirically this yields the highest long-context accuracy among the
methods we evaluate (Section~\ref{sec:results-multipole}).  Its
weakness is that clusters are \emph{data-dependent and recomputed
during generation}: the centroids must be re-estimated as the KV
cache grows, and the keys belonging to a selected cluster are
scattered (non-contiguous) in memory.  The per-step computation
therefore has an irregular memory-access pattern that does not
amortize across decode steps the way a fixed page layout does, and
the decode-time wall clock is substantially slower than the
page-based methods at the same selected-token budget.
This is the source of the accuracy--throughput Pareto trade-off that
Section~\ref{sec:results-multipole} reports in detail.


% =====================================================================
\section{Frequency-Domain Signal Compression}\label{sec:related-freq}
% =====================================================================

The DCT proxy of Section~\ref{sec:proxy} is a low-rank linear
projection of a key sequence onto its leading discrete-cosine
coefficients.  We collect here the small set of background facts that
make this projection well-behaved for our application.

\paragraph{Discrete cosine transform.}
The type-II DCT used in Eq.~\eqref{eq:dct} is an orthonormal linear
transform that asymptotically diagonalizes the second-order
autocorrelation matrix of a stationary Markov-1 signal --- this is
the classical motivation for its dominance in image and audio
codecs (e.g.\ JPEG)~\cite{rao1990dct}.  Orthonormality means that energy is
preserved under DCT, and Parseval's theorem applies coefficient by
coefficient: discarding the higher-frequency bins reduces the
sequence's $\ell_2$ norm by exactly the energy of those bins.
Whether the resulting low-rank reconstruction is a good summary of
the original signal therefore reduces to whether the signal's energy
is in fact concentrated in the leading bins, which is the empirical
question raised in Section~\ref{sec:insight}.

\paragraph{Lowpass-IDCT as a linear projection.}
Concatenating a length-$P$ DCT, a truncation that keeps the first $C$
coefficients, and a length-$C$ inverse DCT yields a single
$C \times P$ matrix that can be precomputed once and applied as a
small dense matrix multiply at runtime.  This is the operator
$M \in \mathbb{R}^{C \times P}$ of Eq.~\eqref{eq:proxy}.  All of the
mass-recall and accuracy results in Chapter~\ref{chap:results} use
this single fixed $M$; nothing about the proxy is learned.

\paragraph{Concurrent frequency-domain work.}
The closest contemporary system is
FreqKV~\cite{kai2026freqkv}, which uses a DCT-based representation
to compress the KV cache for context-window extension rather than
for sparse decoding.  FreqKV operates in the prefill stage and
modifies the cache layout itself, whereas DCT-Page leaves the cache
unchanged and uses DCT only as a scoring proxy during decode.  The
two systems are complementary --- one could in principle combine
FreqKV's frequency-domain prefill compression with DCT-Page's
frequency-domain decode-time scoring --- but in this thesis we
restrict attention to the decode-time use of DCT.
